\section*{ACKNOWLEDGMENT}
Large language model assistance (DeepSeek-V4.1 Flash, a variant within the
DeepSeek-V4 model family \cite{b33}, run in max-reasoning mode via OpenRouter) was
used for language editing throughout the manuscript. The LLM-judge feedback scoring
central to the study's methodology is a described experimental component
(Section~\ref{sec:data}), not AI-generated article text, and is cited at each point
of use in the main text. All conceptual framing, analysis, and interpretation are
the authors' own.

% TODO (Overleaf): add the three author photos. Replace
%   \begin{IEEEbiography}{Name}
% with
%   \begin{IEEEbiography}[{\includegraphics[width=1in,height=1.25in,clip,keepaspectratio]{photo.jpg}}]{Name}
% and upload the photos with the manuscript. The class inserts "PLACE PHOTO HERE"
% when the optional photo argument is absent.
\newpage

\begin{IEEEbiography}{Domen Jer\v{s}ek} is currently completing an undergraduate degree in Applied
Mathematics at the Faculty of Mathematics and Physics, University of Ljubljana,
Slovenia, having enrolled in an additional year of study beyond the standard
program to focus on his research and professional development, with expected
completion in 2027.

Since November 2024, he has worked in a student research capacity with the
Artificial Intelligence Department (E3) at the Jo\v{z}ef Stefan Institute,
Ljubljana, contributing to project work within the European Horizon project
HumAIne. He is a co-author of two prior conference papers presented at the
Information Society (SIKDD) conference. His research interests include
retrieval-augmented generation, large language models as evaluators, and
feedback-driven ranking for enterprise information systems.
\end{IEEEbiography}


\begin{IEEEbiography}{Blaž Kažič} is a researcher at the Artificial Intelligence Laboratory, Jožef Stefan Institute, Slovenia. He received a degree in electrical engineering from the University of Ljubljana and a Ph.D. in Information and Communication Technologies from the Jožef Stefan International Postgraduate School. His research spans machine learning, probabilistic modeling, and time series forecasting, with applications in human mobility, energy consumption, and manufacturing demand. His industry experience includes work on digital twins, manufacturing optimization, and anomaly detection for financial systems. His current work focuses on retrieval-augmented generation, knowledge-grounded language models, and the design and evaluation of agentic systems deployed in organizational workflows. He also investigates the safety and reliability of language models, particularly how changes to refusal behavior affect model capabilities and tool use.
\end{IEEEbiography}



\begin{IEEEbiography}{Klemen Kenda} is a researcher at the Artificial
Intelligence Department (E3), Jožef Stefan Institute. He received his PhD in Information and
Communication Technologies from the Jožef
Stefan International Postgraduate School in 2024,
following a diploma in physics from the University
of Ljubljana. His research focuses on machine
learning for heterogeneous and streaming data,
data fusion, contextual modelling, and evolutionary
approaches to feature selection, with recent work
spanning human-in-the-loop approaches,
knowledge graphs, large language models, and
human-AI collaboration. He has contributed to
over a dozen European research projects since
2011 as a researcher or work package leader
across FP7, Horizon 2020, Horizon Europe, and
Marie Skłodowska-Curie programmes. His work on
human-in-the-loop machine learning has included
applications in smart cities, IT, and healthcare.
Alongside his academic work, he has worked in
several companies in industry in roles including
Lead Research Scientist and Data Scientist.
\end{IEEEbiography}



\EOD

\end{document}
